\section{Vérification}
\label{sec-verification}

\subsection{Cadre}

La campagne en cours adopte les définitions de la norme ASME V\&V~10~\cite{asme2019vv10} et d'Oberkampf et Roy~\cite{oberkampf2010vv}. La vérification établit que le code résout correctement ses équations, contre une solution analytique, une propriété mathématique ou une implémentation indépendante, sans donnée expérimentale. La validation établit que les équations représentent la roche, contre un essai qui n'a pas servi au calage. Un test de non-régression, dont la référence est une sortie antérieure de rockim, ne relève d'aucune des deux : il garantit que le code n'a pas changé, pas qu'il est juste (\cle{docs/VV_campagne.md}, l.~43-56). Cette section rassemble ce qui relève de la vérification ; la validation fait l'objet de la \cref{sec-validation}.

Les chiffres de cette section sont de deux natures. Les uns sont écrits dans les documents du dépôt. Les autres ont été rejoués le 6 octobre 2026 pour ce rapport, sur un exécutable compilé sous Linux à partir du commit \cle{f0209ef}, à un fil ; ils sont signalés comme tels.

\subsection{La suite de non-régression}

Le script \cle{tools/verify_suite.py} lance l'exécutable sur une configuration ou une sous-commande d'auto-test, lit les grandeurs imprimées et les compare à des références par une tolérance absolue propre à chaque contrôle. Il connaît aussi des contrôles qualitatifs : présence du verdict interne de l'auto-test, travail d'un amortisseur négatif ou nul, absence d'un déclenchement. Les références sont figées sur une base Linux du 11 août 2026, et un fichier de références par plateforme existe pour macOS ARM, car les tirages aléatoires de Voronoï et de Weibull diffèrent d'une bibliothèque standard à l'autre à graine égale.

La suite compte 118 tests, comptés en exécutant sa définition : 51 au niveau rapide, 56 de plus au niveau complet, 11 de plus au niveau exhaustif. Plusieurs documents en donnent un autre nombre (12, 44, 48 ou 104) : ces nombres sont datés et périmés, la suite ayant grossi par addition.

Le document de campagne classe ces 118 tests selon la nature de leur référence (\cref{tab-classement}). Seuls les 56 tests des quatre premières catégories relèvent de la vérification. Les 24 tests qui comparent un pic à $f_t$ figent l'écart mesuré un jour donné, de sorte qu'un écart de $1{,}4~\%$ ou de $15~\%$ passerait de la même façon s'il avait été enregistré tel quel. Les 38 tests de régression ne démontrent pas la justesse. Ce classement est déclaré dans le document ; la liste nominative test par test n'est pas outillée.

\begin{table}[htbp]
\centering
\caption{Classement des 118 tests de la suite selon la nature de leur référence (\cle{docs/VV_campagne.md}, Partie I §2).}
\label{tab-classement}
\begin{tabularx}{\linewidth}{@{}X r X@{}}
\toprule
catégorie & tests & référence \\
\midrule
analytique au point matériel ou sans maillage & 11 & formule exacte, VUMAT Fortran, conservation à deux corps \\
analytique sur une structure & 8 & barre en traction, allongement et réaction exacts \\
invariant au repos & 26 & aucune rupture, aucune injection sous charge nulle \\
invariant d'équivalence & 11 & deux réglages qui doivent coïncider, ou un poste qui ferme le bilan \\
pic contre $f_t$, écart figé & 24 & écart à $f_t$ mesuré un jour donné \\
régression pure & 38 & sorties antérieures de rockim \\
\bottomrule
\end{tabularx}
\end{table}

Rejouée pour ce rapport au niveau rapide, la suite passe 50 tests sur 51 en 19,5~min. Le seul échec, \cle{t1_toolcontact_penalty}, porte sur le banc de raclage en pénalité, cas construit pour pomper de l'énergie et déclaré chaotique : il donne un rapport d'injection de $2{,}985$ pour $4{,}426$ attendu, écart déjà consigné pour Linux dans le journal des modifications. La suite doit être lancée depuis la racine du dépôt, faute de quoi des chemins de maillage relatifs ne sont pas trouvés. Sous macOS ARM, le document de campagne déclare 39 tests sur 50, les 11 échecs étant tous des pics à écart figé ou des régressions dont les grandeurs post-pic sont déplacées par l'arithmétique flottante ; aucun test de vérification n'y échoue.

\subsection{Auto-tests}
\label{sec-selftests}

Douze sous-commandes \cle{rockim selftest-*} vérifient une brique sans maillage ou sur un maillage minimal ; huit sont branchées dans la suite. Le \cref{tab-selftests} donne leurs résultats rejoués.

\begin{table}[htbp]
\centering
\caption{Auto-tests rejoués le 6 octobre 2026 sous Linux, à un fil.}
\label{tab-selftests}
\begin{tabularx}{\linewidth}{@{}l X X@{}}
\toprule
auto-test & ce qui est vérifié & résultat rejoué \\
\midrule
\cle{potential2d} & collision élastique sans frottement par le potentiel de Munjiza 2D, frontale et oblique & $\Delta E_c/E_{c0} = 3{,}72\times10^{-12}$ et $5{,}12\times10^{-7}$ ; quantité de mouvement à $4{,}9\times10^{-15}$ \\
\cle{potential3d} & même essai entre tétraèdres & $2{,}05\times10^{-8}$ et $7{,}3\times10^{-9}$ ; quantité de mouvement à $4{,}7\times10^{-15}$ \\
\cle{potvolume3d} & même essai, force par volume de recouvrement & $7{,}1\times10^{-15}$ et $1{,}21\times10^{-11}$ \\
\cle{toolcontact} & algèbre du contact d'outil en vitesse ; chaîne de $N$ masses frappée à 10~m/s & écart relatif maximal $0{,}00996$ ; $N = 400$ : sortie à $19{,}936$~m/s pour 20, perte $0{,}2475~\% \approx 1/N$ \\
\cle{mc} & plateaux plastiques de Mohr-Coulomb contre $\sigma_c = 2c\cos\varphi/(1-\sin\varphi)$ & compression $-3{,}0\times10^{-9}~\%$ ; triaxial au plus $2{,}1\times10^{-9}~\%$ ; traction $-0{,}126~\%$ \\
\cle{pdc3d} & géométrie du couteau PDC chanfreiné, avec variantes qui doivent échouer & 14 contrôles passés, aucun échec \\
\cle{saksala2011}, \cle{dpdfh} & lois au point matériel & s'exécutent ; la comparaison au Fortran n'est pas faite par la sous-commande \\
\bottomrule
\end{tabularx}
\end{table}

Les écarts de $8\times10^{-14}$ et $4{,}7\times10^{-12}$ annoncés pour les portages des VUMAT de Saksala et du DP-DFH ne sont pas rejouables depuis le dépôt seul : la sous-commande écrit la trajectoire de contrainte, et la référence Fortran n'est pas dans le dépôt. Quatre autres auto-tests (empreinte triaxiale, endommagement plastique, fissure fixe, DFH reformulé) ne sont pas dans la suite et n'ont pas été rejoués. En pénalité, le même nœud au repos repart en un pas à 373~m/s, soit $18{,}7$ fois la borne physique $2v$, et un amortissement de Cundall de $0{,}05$ fait perdre $14{,}7~\%$ de la vitesse de sortie à la chaîne de masses.

\subsection{Bilan d'énergie}
\label{sec-B4}

Le bilan d'énergie, ajouté le 14 août 2026 en FDEM 2D et 3D, applique le théorème de l'énergie cinétique aux nœuds : la variation d'énergie cinétique égale la somme des travaux des familles de forces, plus un résidu (\cle{DOCUMENTATION_rockim.md}, l.~2622-2641). Les postes sont les éléments, les joints, les amortisseurs, le contact général et son frottement, l'amortissement de Cundall, les frontières, l'outil, les platines, les forces de volume et un poste d'intégration. Ce dernier corrige exactement l'écart du saute-mouton, $f^2\Delta t^2/(2m)$ par nœud et par pas, entre les compteurs qui lisent la vitesse d'avant le pas et le théorème discret qui veut la moyenne : sans lui, le résidu de la percussion 2D valait $91~\%$, avec lui $0{,}017~\%$. Le bilan est déclaré bon quand le résidu reste sous $1~\%$ du flux brut échangé.

Les résidus mesurés sont de l'ordre de l'arrondi en élastique et petits sur les impacts. Une barre en traction rejouée pour ce rapport donne $-5{,}6\times10^{-17}$~J en FDEM 3D et $3{,}4\times10^{-17}$~J en éléments finis 3D. La documentation donne $0{,}017~\%$ sur la percussion 2D, $0{,}005~\%$ sur la percussion 3D et $-5\times10^{-24}$~J à charge nulle entre deux corps ; le calcul St Anne de $300~\mu$s ferme à $1{,}5\times10^{-7}~\%$.

Le bilan est une identité comptable et ne voit pas une création d'énergie logée dans un poste compté. La coupe PDC 2D du 18 août injectait 408 fois le travail de l'outil pendant que le bilan affichait un résidu de $1{,}9\times10^{-10}~\%$. Pour cette raison, des indicateurs dédiés s'ajoutent au bilan pour le contact d'outil (rapport d'injection, borne $2v$), et la loi de joint se juge sur cycle fermé (\cref{sec-decharge}). Le garde-fou \cle{budgetAbortPct} arrête un calcul dont le résidu dépasse un seuil ; sans compter la pesanteur dans le bilan (\cle{energyBodyForces = on}), il coupait à tort un calcul sain sur un résidu de $175~\%$ qui était le travail de la pesanteur.

Le bilan a une réserve d'instrumentation en FDEM 3D : le poste des éléments est unique, exact en élastique et approché sous loi de volume ou plafond, et la dissipation plastique du volume n'est pas séparée (\cle{AUDIT_loi_adaptative_2026-09-11.md}, §5).

\subsection{Contrôles à charge nulle}

Un contrôle à charge nulle laisse un modèle au repos, ou sous une vitesse de $10^{-12}$~m/s, et exige zéro joint rompu et un travail d'amortisseur négatif ou nul. La suite en compte 26, en 2D et en 3D, un par option de loi ou de contact ajoutée. Ils ont attrapé de vrais défauts : 158 joints rompus à charge nulle sur un maillage de Delaunay avant correction, et 5 joints rompus par des tétraèdres exactement tangents dans le contact par potentiel 3D. Tous ceux du niveau rapide passent, rejoués pour ce rapport.

\subsection{Nombre de fils et déterminisme}
\label{sec-fils}

Les résultats dépendent du nombre de fils. À un fil, les sorties sont identiques au bit à celles du calcul série ; à nombre de fils fixé, elles sont déterministes ; entre deux nombres de fils, l'ordre des sommes change et les pics s'écartent. La documentation annonce environ $2~\%$ ; le cas mesuré le plus documenté donne 7\,916,75~N à 4 fils contre 7\,791,04~N à 2 fils, soit $1{,}6~\%$ (\cle{tools/BITID.md}). Le banc de raclage en pénalité rompt 8 joints au lieu de 2 hors d'un fil. La suite se certifie donc à un fil.

L'identité au bit sert à montrer qu'une option nouvelle ne change rien quand elle est absente. Huit configurations courtes, ancrées à 4 fils le 5 septembre 2026, voient l'empreinte de leurs sorties comparée à chaque modification (\cle{tools/bitid.py}) ; toute ancre changée exige une ligne au journal des modifications. Cette discipline protège contre le changement involontaire ; elle ne prouve pas la justesse.

\subsection{Comparaisons à un autre code et à Hertz}

Une comparaison à Abaqus avec la loi de Saksala sur le même maillage, une sphère à 8~m/s, a été faite le 31 juillet avec des critères fixés avant calcul : pic filtré $+7{,}9~\%$, énergie $+0{,}9~\%$, pénétration $+1{,}1~\%$, restitution $0{,}111$ contre $0{,}059$. Une seconde comparaison avec une loi d'endommagement plastique est passée de $-20~\%$ à $+2{,}8~\%$ après l'ajout de quatre clés choisies pour se rapprocher d'Abaqus, ce que le document qualifie d'ajustement et non de vérification. Les fichiers de ces comparaisons ne sont pas dans le dépôt. La figure \cle{figures_impact/compare_fdelta_abaqus_rockim.png} compare le bruit de deux cas différents et non leurs résultats.

La réponse en force d'un impact élastique a été confrontée à Hertz sur quatre maillages (\cref{fig-hertz}). La réponse globale suit Hertz à $-25~\%$ à $+30~\%$ près en fin de charge, sans converger vers le rapport 1 au raffinement, et le pic local de contrainte n'est pas résolu : il passe de 266 à 2\,931~MPa quand la taille d'élément passe de $4{,}74$ à $0{,}80$~mm, pour une pression de Hertz de 4\,114~MPa. Le document de campagne écrit pourtant que le contact n'a pas de référence analytique en force ; le banc de la sphère sur un plan, préparé, reprend cette question avec des critères écrits (\cref{sec-validation}).

\begin{figure}[htbp]
\centering
\includegraphics[width=\linewidth]{hertz_compare}
\caption{Impact élastique d'une sphère sur quatre maillages de roche : force contre enfoncement, rapport à Hertz et pic de von Mises contre taille d'élément (\cle{figures_impact/plot_hertz.py}).}
\label{fig-hertz}
\end{figure}

\subsection{Audits}

Quatre audits ont précédé la campagne de vérification. La revue de fiabilité 3D du 14 août a établi un plan à critères de réfutation après l'instabilité de $2{,}4$~MJ, dont le garde-fou d'énergie ; le cas explosif n'a jamais été rejoué depuis les correctifs. L'audit de la loi adaptative du 11 septembre a confronté la loi codée à une note de loi constitutive : 2 briques conformes, 11 partielles, aucune absente, et six mécanismes dissipatifs parasites armés par défaut. L'audit du 13 septembre, conduit par quatre agents en lecture seule, a relevé la pénalité de contact absente du budget du pas, la contrainte normale divisée par deux sous la branche parabolique et le contact roche-roche dix fois trop raide ; les corrections sont des options désactivées par défaut. Le plan de robustesse du 5 septembre a diagnostiqué un mode de panne dominant, le défaut silencieux, et a conduit au refus des clés inconnues et à l'ancre d'identité au bit ; les références Linux et Windows de la suite et l'intégration continue qu'il prévoyait ne sont pas faites.
